\section{Harness 的三大组成部分}

Birgitta 将 OpenAI 团队描述的 harness 组件归纳为三个类别，这些组件混合使用了确定性（deterministic）和基于 LLM 的方法。

\subsection{Context Engineering（上下文工程）}

第一个组成部分是\textbf{上下文工程}，即持续增强代码库中的知识库，同时让 agent 能够访问动态上下文信息。

\begin{importantbox}{Context Engineering 的两个维度}
\begin{enumerate}[leftmargin=1.5em]
    \item \textbf{静态知识库}：在代码库中持续维护和增强的文档、设计决策记录、架构约定等
    \item \textbf{动态上下文}：agent 可以实时访问的可观测性数据（observability data）、浏览器导航等运行时信息
\end{enumerate}
上下文工程不仅仅是编写 Markdown 规则文件那么简单------代码设计本身就是上下文的重要组成部分。
\end{importantbox}

这一点非常重要：OpenAI 团队的上下文工程不仅涉及知识库的整理（curation），还涉及大量的\textbf{设计工作}------代码的设计本身就是上下文的核心组成部分。这意味着，好的代码架构不仅对人类开发者有益，对 AI agent 同样至关重要。

\subsection{Architectural Constraints（架构约束）}

第二个组成部分是\textbf{架构约束}，通过确定性和 LLM 双重机制进行监控和执行。

具体包括：
\begin{itemize}[leftmargin=1.5em]
    \item \textbf{基于 LLM 的 agent 监控}：AI agent 在生成代码时自我检查是否符合架构约定
    \item \textbf{确定性自定义 linter}：不依赖 AI 的硬性规则检查工具
    \item \textbf{结构性测试（structural tests）}：验证代码结构是否符合预期的自动化测试
\end{itemize}

\begin{knowledgebox}{确定性 vs LLM 驱动的检查}
在 harness 中，确定性工具（如自定义 linter、结构性测试）和 LLM 驱动的检查各有优势：

\begin{table}[H]
\centering
\begin{tabular}{lll}
\toprule
\textbf{维度} & \textbf{确定性工具} & \textbf{LLM 驱动} \\
\midrule
可靠性 & 100\% 一致 & 可能有误判 \\
灵活性 & 仅处理预定义模式 & 可理解语义意图 \\
维护成本 & 需手动更新规则 & 可自适应新模式 \\
适用场景 & 硬性边界、格式规范 & 设计意图、模式匹配 \\
\bottomrule
\end{tabular}
\end{table}

两者的结合是 harness 设计的关键------确定性工具提供可靠的底线保障，LLM 则覆盖更复杂的语义检查。
\end{knowledgebox}

\subsection{Garbage Collection（``垃圾回收''）}

第三个组成部分是\textbf{周期性运行的清理 agent}，负责发现文档中的不一致性或架构约束的违规行为，对抗代码库的\textbf{熵增}（entropy）和\textbf{衰减}（decay）。

\begin{warningbox}{代码库的熵增问题}
在大规模代码库中，即使每次单独的修改都是合理的，随着时间推移，整体代码质量也会不断下降------这就是软件工程中的``熵增''现象。传统团队依靠代码审查和定期重构来对抗熵增；在 AI 驱动的开发中，这种熵增可能会更快发生，因此需要专门的 ``garbage collection'' agent 来持续清理和修复。
\end{warningbox}

\subsection{迭代反馈循环}

OpenAI 团队特别强调了 harness 的\textbf{迭代性}。他们的核心方法论可以概括为：

\begin{importantbox}{Harness 的迭代改进方法论}
``当 agent 遇到困难时，我们将其视为一个信号：识别缺失的内容------工具、护栏、文档------并将其反馈回代码库，始终由 Codex 本身来编写修复。''

这形成了一个\textbf{正反馈循环}：agent 的失败驱动 harness 的改进，而改进后的 harness 又提升了 agent 的能力。
\end{importantbox}

\subsection{Birgitta 指出的缺失部分}

Birgitta 注意到，OpenAI 文章中描述的所有措施都集中在提升长期的\textbf{内部质量}（internal quality）和\textbf{可维护性}（maintainability），但缺少对\textbf{功能和行为验证}（verification of functionality and behaviour）的讨论。这是一个值得关注的 gap------代码结构良好但功能不正确，同样是不可接受的。

\begin{warningbox}{内部质量 vs 外部质量}
OpenAI 的 harness 聚焦于架构约束、代码一致性等\textbf{内部质量}指标，但没有讨论如何验证代码的\textbf{外部质量}------即功能是否正确、行为是否符合预期。一个完整的 harness 应该同时覆盖这两个维度。缺少功能验证的 harness 就像只检查语法不检查语义的编译器。
\end{warningbox}

\subsection{本章小结}

Harness 由三大组成部分构成：上下文工程提供 AI 所需的知识和信息；架构约束通过确定性和 LLM 双重机制执行规则；``垃圾回收'' agent 持续对抗代码库的熵增。三者的迭代协作形成正反馈循环，但目前在功能验证方面仍有缺口。

\newpage

